% =============================================================================
% Chapter 6: Functional Requirements
% =============================================================================

\chapter{Functional Requirements}
\label{ch:functional}

This chapter consolidates the functional requirements of \itc{}, organized
by topic. Numbering is stable: a deprecated requirement keeps its
\reqid{REQ-XX} identifier forever and is tagged \deprecated{}.

\section{Data Loading: \code{itc.load()}}

\begin{reqbox}{REQ-01}{\stable}
\textbf{Unified loading entry point.}\\
\code{itc.load(source, dims=None, version=None, mode="production",
user=None, comment=None)} is the single entry point for loading data into
\itc{}. The \code{source} argument accepts:
\begin{itemize}[noitemsep,topsep=2pt]
  \item a folder path (loads all compatible files);
  \item a single file path (loads as a one-source dataset);
  \item a dictionary mapping coordinate tuples to file paths (most
    traceable, recommended for production); a tuple position may hold
    the sentinel \code{"*"} to declare the dimension swept within the
    file, its coordinates then being read from the file column of the
    same name (OQ-19);
  \item a NumPy array plus a list of variable names;
  \item a pandas DataFrame.
\end{itemize}
When \code{dims} is omitted or set to \code{None}, the result is in
\emph{datapoint mode}: a single synthetic dimension named
\code{datapoint} indexes rows in acquisition order and every column
becomes a variable. A later call to \code{db.pivot(dims=[...])}
reorganizes the data into a structured VarFrame. String-valued columns
are accepted only as dimension coordinates; loading one as a variable
raises \code{DataError} (OQ-21).
\end{reqbox}

\begin{reqbox}{REQ-02}{\stable}
\textbf{Folder load with filename pattern.}\\
\code{itc.load(folder, pattern=...)} accepts an optional \code{pattern}
argument that filters files by filename, supporting both Unix-style glob
strings (\code{"run\_*.csv"}) and Python regular expressions
(\code{r"run\_m\textbackslash d+\_r\textbackslash d+\_b\textbackslash d+\textbackslash.csv"}).
\end{reqbox}

\begin{reqbox}{REQ-03}{\stable}
\textbf{Explicit dictionary loading.}\\
\code{itc.load(\{coord\_tuple: path, ...\}, dims=[...])} accepts a
dictionary mapping coordinate tuples to file paths. Coordinate tuples
correspond positionally to \code{dims}. Tuple-length mismatch with
\code{dims} raises \code{LoadCoordinateError}.
\end{reqbox}

\begin{reqbox}{REQ-04}{\stable}
\textbf{NumPy array loading.}\\
\code{itc.load(arr, names=[...])} accepts a 2-D NumPy array plus a list
of column names with the same length as the second axis of \code{arr}.
The result is a VarFrame in datapoint mode, ready for \code{db.pivot}.
\end{reqbox}

\begin{reqbox}{REQ-05}{\stable}
\textbf{Pandas DataFrame loading.}\\
\code{itc.load(df)} accepts a pandas DataFrame and returns a VarFrame in
datapoint mode using the column names as variable names. The dependency
on pandas is lazy.
\end{reqbox}

\begin{reqbox}{REQ-06}{\stable}
\textbf{Missing combinations filled with NaN.}\\
Combinations of dimension coordinates not covered by any input file are
filled with \code{np.nan}. \itc{} does not raise an error for sparse test
matrices; the unfilled positions are reported by \code{db.diagnostics()}.
\end{reqbox}

\begin{reqbox}{REQ-07}{\stable}
\textbf{Provenance and history at load time.}\\
At load time \itc{} records: source file paths; SHA-256 of their
concatenated content; ISO 8601 timestamp; user identity
(\code{getpass.getuser()+"@"+socket.gethostname()}, overridable via
\code{itc.set\_user(...)} or \code{user=}); operating mode; \itc{}
version; and an optional version tag. The load operation is the first
entry in the History (index 1).
\end{reqbox}

\section{Operating Modes}

\begin{reqbox}{REQ-08}{\stable}
\textbf{Two operating modes.}\\
\itc{} supports two operating modes, declared at load time via
\code{mode="production"} (default) or \code{mode="draft"}, or globally via
\code{itc.set\_mode(...)}. The mode is stored in Provenance and is
visible in \code{print(db)}, \code{db.summary()}, and \code{db.history()}.
\end{reqbox}

\begin{reqbox}{REQ-09}{\stable}
\textbf{Production mode: mandatory history.}\\
In production mode, every operation records itself in History
automatically. No additional argument is required to enable recording.
This is the only acceptable mode for analyses producing official results.
\end{reqbox}

\begin{reqbox}{REQ-10}{\stable}
\textbf{Draft mode: opt-in history.}\\
In draft mode, history entries are recorded only when the user explicitly
passes \code{history=True} to an operation. VarFrame objects created in
draft mode carry \code{mode="draft"} in their Provenance and are clearly
identified by \code{print(db)}. Draft mode is intended exclusively for
initial data exploration.
\end{reqbox}

\begin{reqbox}{REQ-11}{\stable}
\textbf{Draft mode: save guard.}\\
Calling \code{db.save(...)} or any export method on a VarFrame in draft
mode raises \code{DraftModeExportError} by default. The guard is
overridable by passing \code{allow\_draft=True}, which forces a deliberate
second decision and embeds a prominent warning in the output file
metadata. The guard covers result exports (\code{save} and the
\code{to\_*} family); inspection artifacts (\code{db.manifest} files,
\code{db.diagnostics(log=...)} logs) are exempt because they exist to
understand draft data (OQ-22).
\end{reqbox}

\begin{reqbox}{REQ-12}{\stable}
\textbf{Strict mode mixing in binary operations.}\\
When an operation combines two VarFrame objects (\code{db.combine},
\code{pproc.compare}, etc.), both inputs must share the same operating
mode. Mixed-mode combinations raise \code{OperatingModeMixError}; the user
must call \code{db.promote(mode="production")} or
\code{db.demote(mode="draft")} explicitly before the operation. There is
no implicit promotion or demotion.
\end{reqbox}

\begin{noteinfo}[Draft mode is for understanding data, not for results]
Any analysis whose numbers feed a report, a publication, or an engineering
decision must run in production mode.
\end{noteinfo}

\section{Inspection and Diagnostics}

\begin{reqbox}{REQ-13}{\stable}
\textbf{Inspect a datapoint-mode VarFrame.}\\
\code{db.inspect(threshold=20)} analyzes each variable in a datapoint-mode
VarFrame and reports: unique value count, dimension-vs-variable candidacy,
coverage estimate, and ambiguous variables. Output is printed to the
terminal. \code{db.inspect()} is a no-op when the VarFrame is already
structured (multiple dimensions present).
\end{reqbox}

\begin{reqbox}{REQ-14}{\stable}
\textbf{Pivot to a structured VarFrame.}\\
\code{db.pivot(dims=[...], auto\_detect=False)} reorganizes a datapoint-mode
VarFrame into a structured one. With \code{auto\_detect=True}, additional
dimension candidates are detected automatically by inspecting unique-value
counts. Calling \code{db.pivot()} without \code{dims} on an already-structured
VarFrame raises \code{PivotError}.

When two or more datapoints share identical coordinates on every requested
dimension, \code{db.pivot} raises \code{PivotDuplicateError}, listing the
colliding coordinates and suggesting either adding a \code{repeat}
dimension or deduplicating upstream. Real acquisition tables contain
untagged repeats as the norm; silently keeping the first or the last row
would destroy evidence.
\end{reqbox}

\begin{reqbox}{REQ-15}{\stable}
\textbf{Manifest export.}\\
\code{db.manifest(path)} exports a tabular file (CSV or JSON) mapping each
source file to its dimension coordinate values. When a dimension is swept
\emph{within} a single source file (i.e.\ multiple values of that dimension
appear in one file), the manifest records the symbol \code{*} in that
column for that file, indicating ``swept here, see file''. The manifest is
the audit checkpoint between data loading and any further manipulation.
\end{reqbox}

\begin{reqbox}{REQ-16}{\stable}
\textbf{Summary: short and practical.}\\
\code{db.summary()} returns a concise object that prints to the terminal a
one-screen summary: dimensions and their cardinalities; variable list;
min/max/mean per variable (where finite); RAM footprint of the data and
metadata; operating mode; current history index. Returns a \code{Summary}
object with matching attributes. Designed to be the first thing the user
calls after \code{itc.load(...)}.
\end{reqbox}

\begin{reqbox}{REQ-17}{\stable}
\textbf{Diagnostics: print and return.}\\
\code{db.diagnostics(log=None)} prints a structured report to the terminal
and returns a \code{DiagnosticsReport} object with attributes
\code{.missing}, \code{.non\_finite}, \code{.summary}, \code{.partial\_vars},
\code{.coverage}, \code{.n\_missing}, \code{.warnings}, \code{.to\_csv()},
\code{.to\_json()}. The optional \code{log=path} argument writes the
printed output to a \code{.log} file in addition to the terminal.
\end{reqbox}

\section{Core Operations}

\begin{reqbox}{REQ-18}{\stable}
\textbf{Immutable operations.}\\
All operations on a VarFrame return a new VarFrame. No operation modifies
the original in place. Every returned VarFrame carries the Provenance of
the original, an updated History that includes the operation and its
arguments, and an updated state hash.
\end{reqbox}

\begin{reqbox}{REQ-102}{\stable}
\textbf{Array-level immutability.}\\
Structural immutability (DD-03) extends to the arrays themselves: every
NumPy array stored in a VarFrame, UncFrame, or HistoryFrame is set
read-only (\code{writeable=False}) at construction, and \code{to\_numpy}
returns read-only views by default (\code{copy=True} for a writable
copy). A frozen dataclass alone does not prevent
\code{db.vars["CT"].values[0] = 99}; the read-only flag does. Tests must
verify that in-place mutation attempts raise.
\end{reqbox}

\begin{reqbox}{REQ-103}{\stable}
\textbf{State-hash scope and reproducibility contract.}\\
The \code{state\_hash} covers dimension names and coordinates, variable
names and values, the ordered operation sequence with normalized arguments
and comments, and the uncertainty, correlation, and origin-tag arrays when
present. It excludes timestamps, user identity, source paths, and the
\itc{} version (Section~\ref{sec:state-hash}). Loading the same bytes and
applying the same operations on the same platform and NumPy version yields
identical hashes; this is a tested invariant. Cross-platform comparison
uses \code{pproc.compare} with tolerances, never hash equality.
\end{reqbox}

\begin{reqbox}{REQ-19}{\stable}
\textbf{Optional history comment on every operation.}\\
Every operation accepts an optional keyword argument \code{comment=str}
which is stored alongside the History entry, allowing the user to attach
context (e.g.\ \code{"removed bad point at alpha=8.0 from drift in run 42"}).
The comment is preserved in \code{db.history()}, in pipeline exports, and
in \code{.itc} archives.
\end{reqbox}

\begin{reqbox}{REQ-20}{\stable}
\textbf{Selection with operators and frame target.}\\
\code{db.select(filters, Frame="VarFrame")} restricts a VarFrame to the
subset of coordinates matching the given filters. Filter keys are
\code{dim\_op} strings of the form dimension name suffixed by an optional
operator (\code{=} default, \code{>}, \code{>=}, \code{<}, \code{<=},
\code{!=}). Values are scalars or lists.

The optional \code{Frame=} keyword targets the comparison source:
\begin{itemize}[noitemsep,topsep=2pt]
  \item \code{Frame="VarFrame"} (default): filter against the variable
    values themselves;
  \item \code{Frame="HistoryFrame"}: filter against the origin tags,
    e.g.\ \code{db.select(\{"CT": 0\}, Frame="HistoryFrame")} retains only
    points where \code{CT} is original;
  \item \code{Frame="UncFrame"}: filter against the standard uncertainty
    of a variable, e.g.\ \code{db.select(\{"CT<": 0.005\}, Frame="UncFrame")}.
\end{itemize}

Dimension-keyed filters subset coordinates and reduce the grid.
Variable-keyed filters, including every \code{Frame="UncFrame"} and
\code{Frame="HistoryFrame"} filter, act on the interior of the grid, where
the matching cells do not in general form a Cartesian sub-grid. They
therefore mask: the result keeps the full dimensions, non-matching cells
are set to \code{np.nan}, and origin tags are preserved. Coordinates whose
entire slice is masked out are dropped. The History entry records the
filter and the number of masked cells.
\end{reqbox}

\begin{reqbox}{REQ-21}{\stable}
\textbf{At a single coordinate.}\\
\code{db.at(dim=value)} returns a slice at the specified coordinate. The
specified dimension is removed from the result. Equivalent to
\code{db.select(\{dim: [value]\}).squeeze()}.
\end{reqbox}

\begin{reqbox}{REQ-22}{\stable}
\textbf{Squeeze unit-cardinality dimensions.}\\
\code{db.squeeze(along=None)} removes dimensions whose cardinality is one.
If \code{along} is specified, only the named dimension is squeezed (raises
if its cardinality is not 1). If all dimensions become unit-cardinality,
the result has a single \code{datapoint} dimension with one entry, holding
the scalar values.
\end{reqbox}

\begin{reqbox}{REQ-23}{\stable}
\textbf{Expand: add a new dimension.}\\
\code{db.expand(dim\_name, values, axis=None)} adds a new dimension to a
VarFrame, broadcasting the existing arrays across the new axis. The new
axis position defaults to the last position. Memory cost is documented in
the operation's docstring.
\end{reqbox}

\begin{reqbox}{REQ-24}{\stable}
\textbf{Concatenate VarFrames along a shared dimension.}\\
\code{itc.concat(list, along=dim)} concatenates a list of VarFrame objects
along the given dimension. All inputs must share all other dimensions
identically (same coordinates, same units); the values along \code{along}
must be unique across inputs (no overlap, raises
\code{ConcatOverlapError}). Cardinalities along \code{along} may differ
across inputs.
\end{reqbox}

\begin{reqbox}{REQ-25}{\stable}
\textbf{Interpolate: densify or translate axis.}\\
\code{db.interpolate(\{dim: new\_coords\}, method="linear", deg=None,
override=False)} interpolates all variables onto a new grid along the
specified dimension(s). Methods: \code{"linear"}, \code{"cubic"},
\code{"nearest"}, \code{"polyfit"} (with \code{deg=}).

By default, when a value in \code{new\_coords} already exists in the
original dimension, the existing value is preserved (the operation skips
recomputation at that coordinate). Pass \code{override=True} to force
recomputation at all target coordinates.

The HistoryFrame is updated: \code{+1} for points within the convex hull of
the original axis, \code{-1} outside.

Axis translation (replacing a dimension with a variable as the sweep axis) is a special form:
\begin{lstlisting}
db_cl = db.interpolate(
    axisTranslation={"from": "alpha", "to": "CL"},
    method="linear",
)
\end{lstlisting}
The target variable must be monotonic along the source dimension;
otherwise \code{AxisTranslationError} is raised before any computation.
Numerical operations on string-valued dimensions raise
\code{NonNumericDimensionError}.
\end{reqbox}

\begin{reqbox}{REQ-26}{\stable}
\textbf{Fill missing values.}\\
\code{db.fill(along=dim, method="linear", deg=None, window=None,
global\_fit=False)} fills \code{np.nan} entries along the specified
dimension. Methods:
\begin{itemize}[noitemsep,topsep=2pt]
  \item \code{"linear"}: linear interpolation between neighbors;
  \item \code{"nearest"}: nearest-neighbor copy;
  \item \code{"polyfit"}: polynomial fit of degree \code{deg} on a
    moving window of \code{window} points (default behavior, mirroring
    \code{diff}); when \code{global\_fit=True}, a single polynomial of
    degree \code{deg} is fit to the full dimension instead.
\end{itemize}
Filled values are tagged \code{+1} in the HistoryFrame. For
signature consistency with the M1 kernel operations, \code{method}
is being moved to keyword-only: passing it positionally is deprecated
and warns from v0.2.0 (M1 Phase B1 decision).
\end{reqbox}

\begin{reqbox}{REQ-27}{\stable}
\textbf{Average along a dimension.}\\
\code{db.average(along=dim)} collapses one or more dimensions by arithmetic
mean of non-NaN values. Multiple dimensions may be specified as a list.
Reduces VarFrame order. HistoryFrame tags follow the worst-case rule.
\end{reqbox}

\begin{reqbox}{REQ-28}{\stable}
\textbf{Numerical integration.}\\
\code{db.integrate(var, over=[dims], coords=None)} integrates one or more
dimensions numerically. \code{coords="polar"} applies the polar area
element $r\,\mathrm{d}r\,\mathrm{d}\theta$. The result is a VarFrame of
reduced order containing the integrated quantity.

Missing values: by default any \code{np.nan} inside the integration domain
makes the integrated result \code{np.nan} (\code{skipna=False}). This is
the fail-loud default, because a partial quadrature is silently biased.
Passing \code{skipna=True} integrates over populated cells only, records
the choice and the populated fraction in History, and tags the result
\code{+1}.
\end{reqbox}

\begin{reqbox}{REQ-29}{\stable}
\textbf{Smoothing.}\\
\code{db.smooth(along=dim, method=..., **kwargs)} smooths all variables
along the specified dimension. Methods: \code{"savgol"} (Savitzky-Golay;
requires \code{window}, \code{polyorder}), \code{"spline"} (requires
\code{smoothing}), \code{"moving\_avg"} (requires \code{window}). Smoothed
values are tagged \code{+1} in the HistoryFrame.
\end{reqbox}

\begin{reqbox}{REQ-30}{\stable}
\textbf{Numerical differentiation.}\\
\code{db.diff(along=dim, window=5, deg=2, nan\_edges=False)} computes the
numerical derivative of all variables along the specified dimension using
a moving polynomial fit:
\begin{itemize}[noitemsep,topsep=2pt]
  \item a window of \code{window} points centered on each point is fitted
    with a polynomial of degree \code{deg};
  \item the constraint \code{window > deg} is enforced before any
    computation; violation raises \code{FitDegreeError} (the shared
    leaf for the too-few-points-for-degree invariant across
    \code{diff}, \code{smooth}, \code{interpolate}, and
    \code{fitmodel}; unified from the former \code{DiffWindowError} at
    the M1 Phase B1 checkpoint);
  \item the analytical derivative of the fitted polynomial is evaluated at
    the central point;
  \item at boundaries where a centered window is unavailable, the window
    is asymmetric (fewer points on one side), preserving output shape;
  \item if \code{nan\_edges=True}, asymmetric-window points are set to
    \code{np.nan}.
\end{itemize}
Result is stored in \code{db.d[dim]} as a VarFrame whose variables follow
the naming convention \code{dVAR\_dalpha}, \code{dVAR\_dmach}, etc.
\end{reqbox}

\begin{reqbox}{REQ-31}{\stable}
\textbf{Polynomial fit model.}\\
\code{db.fitmodel(along=dim, deg=N, comment=None)} fits a polynomial of
degree $N$ to every variable along the specified dimension and replaces
that dimension with $N+1$ coordinate entries representing the polynomial
coefficients. The new dimension keeps the original name suffixed by
\code{\_coef} as a name and uses coordinate labels of the form
\code{alpha\^{}0}, \code{alpha\^{}1}, $\dots$, \code{alpha\^{}N}, mirroring
the variable's exponent in the polynomial fit. The propagation of
uncertainty into coefficient space is not yet frozen (OQ-24, tied to
OQ-18); until it is, \code{fitmodel} raises when uncertainty is
present rather than guessing (DD-18).
\end{reqbox}

\begin{reqbox}{REQ-32}{\stable}
\textbf{Polynomial fit evaluation.}\\
\code{db.fitvalue(coef\_dims=[...], at=\{dim: array\})} evaluates a
VarFrame produced by \code{fitmodel} at user-specified coordinates,
returning a new VarFrame where the coefficient dimensions are replaced by
the evaluation dimensions. Useful for densifying a fit beyond the original
sweep range. Values within the original fit range are tagged \code{+1} in
the HistoryFrame; beyond the range, \code{-1}. The forward evaluation is
exact and linear in the coefficients, but because a \code{fitmodel}
output carries no coefficient uncertainty until OQ-24 is frozen,
\code{fitvalue} defers with \code{fitmodel} and raises when
uncertainty is present, keeping the coefficient-space treatment
coherent (DD-18).
\end{reqbox}

\begin{reqbox}{REQ-33}{\stable}
\textbf{Compute: string equation derivation.}\\
\code{db.compute("VAR = expression", debug=False, where=None, fill=np.nan,
method="symbolic", comment=None)} derives a new variable from a string
equation, parsed as RPN, evaluated vectorially over the entire VarFrame,
and stored as a new variable.

When at least one variable in the right-hand side has an assigned
uncertainty, the propagation rule of Equation~\eqref{eq:lpu} is applied
automatically using partial derivatives obtained symbolically via chain
rule on the RPN expression tree. Correlation terms are included when a
correlation matrix has been declared via \code{db.set\_correlation(...)}.

The default \code{method="symbolic"} uses the GUM LPU described above.
Passing \code{method="mcm"} switches to Monte Carlo propagation
(Section~\ref{sec:uq}); this is the recommended choice for expressions
containing discrete branches.
\end{reqbox}

\begin{reqbox}{REQ-34}{\stable}
\textbf{Compute: debug mode.}\\
\code{db.compute("VAR = expr", debug=True)} prints a structured debug
report before applying the equation: parsed RPN token list, identified
variables and constants, sample numerical evaluation at one representative
point with all intermediate values, and (when an UncFrame is present) the
partial derivatives, the correlation contributions, and the resulting
uncertainty at that point. Designed to help validate equations before
committing to production runs.
\end{reqbox}

\begin{reqbox}{REQ-35}{\stable}
\textbf{Compute: conditional filter.}\\
\code{db.compute("VAR = expr", where="condition", fill=value)} evaluates
the equation only at points where \code{condition} is satisfied. Behavior
at points where the condition is false:
\begin{itemize}[noitemsep,topsep=2pt]
  \item \code{fill=np.nan} (default): those points become NaN;
  \item \code{fill=<scalar>}: those points get the provided scalar;
  \item \code{fill=None}: if \code{VAR} already exists in the
    VarFrame, those points retain their existing value; if \code{VAR}
    does not exist, those points become NaN.
\end{itemize}
The condition string supports comparisons (\code{>}, \code{<}, \code{>=},
\code{<=}, \code{==}, \code{!=}) and logical operators (\code{and},
\code{or}, \code{not}) referencing any variable present in the VarFrame.
Uncertainty is propagated only for filtered-in points.
\end{reqbox}

\begin{reqbox}{REQ-36}{\stable}
\textbf{Compute: NumPy guard.}\\
When \code{compute} is invoked with no variables carrying uncertainty in
the expression, every NumPy function is allowed (\code{np.round},
\code{np.clip}, \code{np.where}, etc.).

When at least one variable in the expression has assigned uncertainty:
\begin{itemize}[noitemsep,topsep=2pt]
  \item differentiable NumPy functions (\code{np.sin}, \code{np.cos},
    \code{np.sqrt}, \code{np.log}, \code{np.exp}) are silently normalized
    to their RPN equivalents and propagated correctly;
  \item non-differentiable NumPy functions (\code{np.round}, \code{np.clip},
    \code{np.where}) raise \code{UncertaintyCompatibilityError} with the
    suggestion to either compute the derived variable before assigning
    uncertainty, or to switch to \code{method="mcm"}.
\end{itemize}
The guard is per-variable: if a \code{compute} expression references only
variables without assigned uncertainty, no guard is applied even when
other variables in the same VarFrame carry uncertainty.
\end{reqbox}

\begin{reqbox}{REQ-37}{\stable}
\textbf{Combine VarFrames explicitly.}\\
\code{db.combine(other, op, comment=None)} returns a new VarFrame whose
variables are the elementwise combination of the inputs under the named
operation. Supported \code{op} values: \code{"sum"}, \code{"diff"},
\code{"product"}, \code{"ratio"}, \code{"mean"}, \code{"weighted\_mean"}
(requires \code{weights=}). Both inputs must share dimensions and
coordinates. Uncertainty is propagated under the declared correlation
structure of each input plus an optional cross-input correlation passed
via \code{cross\_correlation=}. The operation appears in History as
\code{combine(op="...", with=db\_other.id)}. Python operators
(\code{+}, \code{-}, \code{*}, \code{/}) on VarFrames are intentionally
\emph{not} supported (NREQ-08).
\end{reqbox}

\begin{reqbox}{REQ-38}{\stable}
\textbf{Rotate to a registered axis.}\\
\code{db.rotate(target\_axis, vector\_groups=None, comment=None)} returns
a new VarFrame with detected vector quantities expressed in the target
axis frame. Vector groups are identified by a default naming convention
(\code{(FX, FY, FZ)}, \code{(MX, MY, MZ)}) or declared explicitly via
\code{db.declare\_vector(name, components=[...])}. The rotation matrix
is applied to each detected group; uncertainty is propagated using the
rotation matrix as Jacobian, including covariance if a correlation matrix
is present. Condition-dependent frames
(Section~\ref{sec:condition-frames}) are evaluated per grid point.
Unknown target axes raise \code{AxisNotFoundError};
unresolvable vector groups raise \code{VectorGroupError}.
\end{reqbox}

\begin{reqbox}{REQ-100}{\stable}
\textbf{Moment reference-point transfer.}\\
\code{db.translate\_moments(to\_point, from\_point=None, frame=None,
comment=None)} transfers every declared moment vector group between
reference points via
$\mathbf{M}' = \mathbf{M} + \mathbf{r} \times \mathbf{F}$
(Section~\ref{sec:moment-transfer}). The offset $\mathbf{r}$ is expressed
in the stated frame (default: the current frame). The transfer requires a
resolvable force vector group; otherwise \code{VectorGroupError} is
raised. The Jacobian is exact; covariance between force and moment
channels is included when declared. History records source point, target
point, and frame.
\end{reqbox}

\begin{reqbox}{REQ-101}{\draft}
\textbf{Condition-dependent axes.}\\
An \code{Axis} may be defined either by a constant 3$\times$3 rotation
matrix or parametrically via \code{angles\_from=("alpha", "beta")} with a
named angle convention. Parametric frames are evaluated per grid point
from the referenced dimensions or variables. The built-in \code{"wind"}
and \code{"stability"} frames follow AIAA R-004A-1992~\citep{aiaaR004A}.
When a referenced angle carries uncertainty, its chain-rule sensitivity
$\partial R / \partial \alpha$ enters the propagation. Exactly one of
\code{rotation\_matrix} and \code{angles\_from} must be provided;
providing both or neither raises \code{RotationMatrixError}.
\end{reqbox}

\begin{reqbox}{REQ-107}{\draft}
\textbf{Per-group source frames and multi-frame data.}\\
Loaded data may already carry vector quantities in more than one
coordinate frame (for example rig-frame balance channels alongside
tunnel-frame reference vectors). \code{db.declare\_vector(name,
components=[...], frame=...)} therefore binds each vector group to the
frame its components are currently expressed in; the frame defaults to
the canonical body axis when omitted. \code{db.rotate(target\_axis)}
transforms each declared group from its own source frame to the
target, composing through the canonical body axis, so groups in
different source frames are handled correctly in a single call. A
group whose source frame is not registered raises
\code{AxisNotFoundError}. The source frame of every group is recorded
in History and is part of the VarFrame state hash (REQ-103).
\end{reqbox}

\section{Uncertainty Quantification}
\label{sec:uq}

\begin{reqbox}{REQ-39}{\stable}
\textbf{Assign uncertainties.}\\
\code{db.set\_uncertainty(\{var: value, ...\})} assigns standard
uncertainties (GUM clause~2.3.1) to measured variables. Values may be
absolute (float) or relative (string ending in \code{"\%"}). This call
creates the UncFrame on first use and updates it on subsequent calls.
Uncertainty is per-variable; keys not matching any variable raise
\code{UncertaintyKeyError}.
\end{reqbox}

\begin{reqbox}{REQ-99}{\stable}
\textbf{Two uncertainty components.}\\
\code{db.set\_uncertainty(\{...\}, component="systematic")} (default) and
\code{component="random"} assign the two UncFrame components of
Section~\ref{sec:two-components}. The systematic component is treated as
fully correlated across the points of a variable; the random component as
independent between points. Reductions apply their weights to both
components, which yields the $1/\sqrt{N}$ gain for the random component
only. \code{pproc.statistics} additionally estimates a random component
from the scatter over the collapsed dimension and reports it separately
from the propagated systematic component, following AIAA
S-071A-1999~\citep{aiaaS071A}. Every report and export states components
separately and combined ($u = (u_{\mathrm{sys}}^2 +
u_{\mathrm{rand}}^2)^{1/2}$).
\end{reqbox}

\begin{reqbox}{REQ-40}{\stable}
\textbf{Declare correlations.}\\
\code{db.set\_correlation(\{(var\_a, var\_b): r\_ab, ...\})} declares
correlation coefficients between pairs of variables. Default is full
independence ($r=0$ for all pairs). The correlation matrix is stored
alongside the UncFrame and is consulted by every uncertainty
propagation operation. Symmetry and the constraint $|r| \le 1$ are
validated; violations raise \code{CorrelationMatrixError}. Variables
referenced but not present in the VarFrame raise
\code{CorrelationKeyError}.
\end{reqbox}

\begin{reqbox}{REQ-41}{\stable}
\textbf{Symbolic propagation via RPN chain rule.}\\
The default propagation engine implements Equation~\eqref{eq:lpu}
(GUM clause~5, Law of Propagation of Uncertainty) using partial
derivatives computed symbolically via chain rule on the RPN expression
tree. The engine handles arbitrary continuous nonlinearities; the
linearization assumption is the GUM standard one and is documented in the
debug output (REQ-34).
\end{reqbox}

\begin{reqbox}{REQ-42}{\stable}
\textbf{Monte Carlo propagation for discrete-branch models.}\\
When an expression contains a discrete branch (conditional logic that
selects different equations based on input values), symbolic propagation
is unsound. \code{db.compute(..., method="mcm", n\_samples=10000)} switches
to Monte Carlo propagation following GUM Supplement~1~\citep{gum2008s1},
sampling each input from its declared distribution (Gaussian by default,
correlated according to \code{set\_correlation}). The output uncertainty
is the sample standard deviation. Continuous nonlinearities do not require
this method: the symbolic propagation handles them correctly.
\end{reqbox}

\begin{reqbox}{REQ-43}{\stable}
\textbf{Sensitivity analysis as native capability.}\\
The combination of \code{set\_uncertainty} and \code{set\_correlation}
on input variables (representing manufacturing tolerances, measurement
errors, or model uncertainties) with \code{compute} or \code{itc.aerospace}
computations natively yields a sensitivity analysis: the propagated
uncertainty on the output represents the expected variation due to the
specified input variability. No separate sensitivity module is required
(P-12, DD-09).
\end{reqbox}

\begin{reqbox}{REQ-98}{\stable}
\textbf{Uncertainty propagation through every operation.}\\
Every operation declares its effect on the UncFrame (DD-18).
Table~\ref{tab:unc-ops} is normative. The weight-based rows (subset,
interpolate, fill, and the reductions) are linear in the variable
values, so propagation through the operation weights is exact; the
Jacobian rows (\code{rotate}, \code{translate\_moments},
\code{combine}, \code{compute}) linearize a nonlinear map, and the
provisional rows (\code{smooth}, \code{diff}, \code{fitmodel},
\code{fitvalue}) raise until their rule is frozen. No operation
silently drops or silently carries an UncFrame. The \code{smooth},
\code{diff}, and \code{fitmodel} rows remain provisional until OQ-18
and OQ-24 are resolved during v0.2.0 implementation; an operation
whose weight rule is not yet frozen raises when uncertainty is
present rather than guessing. The M0 rows were validated against the
implementation on 2026-07-21; the M1 weight-based rows
(\code{interpolate} and the reductions) were validated at the M1
Phase B1 checkpoint, with the mixed-sign systematic behavior of the
cubic and polyfit kernels carried as a validation item alongside
OQ-18.

\begin{table}[H]
\centering
\caption{Normative UncFrame semantics per operation.}
\label{tab:unc-ops}
\small
\begin{tabularx}{\textwidth}{@{}lX@{}}
\toprule
\textbf{Operation} & \textbf{UncFrame effect}\\
\midrule
\code{select}, \code{at}, \code{squeeze}, \code{concat}, \code{expand},
\code{pivot} & Components are subset, concatenated, or broadcast
unchanged.\\
\code{interpolate}, \code{fill} & Propagated through the
interpolation or fit weights, per point, both components.\\
\code{average}, \code{integrate}, \code{statistics} & Propagated through
the reduction weights: random component with independence between points
(the $1/\sqrt{N}$ gain), systematic component as fully correlated (no
gain).\\
\code{smooth}, \code{diff} & Propagated through the kernel or moving-fit
weights, same component rules as reductions.\\
\code{fitmodel}, \code{fitvalue} & Coefficient-space rule not yet
frozen (OQ-24); both raise when uncertainty is present until
validated. \code{fitvalue}'s forward evaluation is exact and linear,
but it defers with \code{fitmodel} so the coefficient story stays
coherent.\\
\code{rotate}, \code{translate\_moments}, \code{combine}, \code{compute} &
Exact Jacobian propagation with covariance (REQ-37, REQ-38, REQ-41,
REQ-100).\\
\bottomrule
\end{tabularx}
\end{table}
\end{reqbox}

\section{Expression Engine}
\label{sec:rpn-engine}

\begin{reqbox}{REQ-44}{\stable}
\textbf{Expression operator coverage.}\\
Expressions are parsed with the standard-library \code{ast} module and
compiled to the \itc{} operator tree (DD-20). The engine supports:
\begin{itemize}[noitemsep,topsep=2pt]
  \item arithmetic: \code{+}, \code{-}, \code{*}, \code{/}, \code{**};
  \item trigonometric: \code{sin}, \code{cos}, \code{tan}, \code{asin},
    \code{acos}, \code{atan}, \code{atan2};
  \item miscellaneous: \code{sqrt}, \code{abs}, \code{log}, \code{log10},
    \code{exp};
  \item constants: \code{pi}, \code{e}.
\end{itemize}
Each operator implements \code{evaluate()}, \code{d\_da()}, and (where
applicable) \code{d\_db()}, and is independently testable. Property-based
tests are required (Chapter~\ref{ch:nonfunctional}).
\end{reqbox}

\section{Processors: \code{itc.pproc()}}

\begin{reqbox}{REQ-45}{\stable}
\textbf{Processor protocol.}\\
Processors are objects satisfying the \code{Processor} typing protocol:
\begin{lstlisting}
class Processor(Protocol):
    name: str
    version: str
    def info(self) -> None: ...
    def validate(self, db: VarFrame) -> None: ...
    def __call__(
        self,
        db: VarFrame,
        *,
        report: str | None = None,
        comment: str | None = None,
        force: bool = False,
    ) -> VarFrame: ...
\end{lstlisting}
The canonical lifecycle is \code{validate} $\to$ \code{info} $\to$
\code{apply} $\to$ \code{report}.
\end{reqbox}

\begin{reqbox}{REQ-46}{\stable}
\textbf{Processor factory.}\\
\code{itc.pproc(name\_or\_path, config=\{\})} creates a processor either by
registered name (e.g.\ \code{"WT\_propeller"}) or by path to a
\code{.itceq} file. Configuration keys override defaults declared in
\code{[constants]}. Unknown processor names raise
\code{ProcessorNotFoundError} listing the registered alternatives.
\end{reqbox}

\begin{reqbox}{REQ-47}{\stable}
\textbf{Processor idempotence policy.}\\
A processor may declare \code{idempotent: bool} as a class attribute
indicating whether reapplication is meaningful. By default
(\code{idempotent=False}), \code{pproc(pproc(db))} raises
\code{ProcessorIdempotenceWarning} and refuses to re-run unless the
caller passes \code{force=True}. Reapplication, when allowed, always
emits the warning and records the second application in History as a
distinct entry. There is no silent no-op.
\end{reqbox}

\begin{reqbox}{REQ-48}{\stable}
\textbf{\code{.itceq} file format.}\\
Processors are loadable from \code{.itceq} files
(Section~\ref{sec:itceq}). An \code{.itceq} file fully defines a
reproducible workflow and is version-controlled alongside the data. The
parser detects cyclic equations and raises \code{ItceqCycleError} at parse
time, before any computation. By default, equations are evaluated in file
order; passing \code{auto\_sort=True} to the parser enables topological
sorting by dependency, in which case the parser reports the resolved
order to the user as feedback.
\end{reqbox}

\begin{reqbox}{REQ-49}{\stable}
\textbf{Statistics over a repeat dimension.}\\
\code{pproc.statistics(db, along="repeat")} collapses a dimension (by convention named \code{repeat} but any dimension is acceptable) into:
\code{VAR\_mean}, \code{VAR\_std}, \code{VAR\_ci95}, \code{VAR\_min},
\code{VAR\_max}, \code{VAR\_N}. The collapsed dimension is absent from
the output VarFrame.
\end{reqbox}

\begin{reqbox}{REQ-50}{\stable}
\textbf{N-way dataset comparison with named inputs.}\\
\code{pproc.compare(reference=db\_ref, **named\_inputs)} compares an
arbitrary number of VarFrame objects against a reference. Inputs are
keyword-only and must be named (e.g.\ \code{case\_a=db\_a},
\code{case\_b=db\_b}). The output is a VarFrame containing, for each
input variable and each named non-reference input \code{name\_i}:
\code{VAR\_ref}, \code{VAR\_name\_i}, \code{VAR\_delta\_name\_i},
\code{VAR\_delta\_pct\_name\_i}, and (when both inputs have UncFrame)
\code{VAR\_delta\_sigma\_name\_i}. Mismatched grids without \code{reference}
raise \code{GridMismatchError}; with \code{reference}, all inputs are
interpolated onto the reference grid via \code{db.interpolate}.
\end{reqbox}

\begin{reqbox}{REQ-51}{\stable}
\textbf{LaTeX-backed PDF report generation.}\\
\code{pproc.report(db, output=path, title=..., plots=True, tables=True,
warnings=True, keep\_tex=False)} generates a structured PDF using a
LaTeX backend. The PDF contains processor metadata, variable summary,
diagnostic plots, tabulated results, and warnings. Equivalently invokable
via \code{pproc(db, report=path)}. When \code{keep\_tex=True} is passed,
the underlying \code{.tex} sources and any auxiliary files are retained
alongside the output PDF. The dependency on the LaTeX toolchain
(\code{tectonic} preferred, \code{pdflatex} fallback) is checked lazily.
\end{reqbox}

\begin{reqbox}{REQ-52}{\stable}
\textbf{Aerodynamic polar extraction.}\\
The built-in processor \code{"aero\_polar"} extracts classical aerodynamic
polar parameters from a VarFrame that uses $C_L$ as its primary axis (via
\code{interpolate} with \code{axisTranslation} from $\alpha$ if needed):

\begin{tabularx}{\textwidth}{@{}llX@{}}
\toprule
\textbf{Output}      & \textbf{Required} & \textbf{Description}\\
\midrule
\code{CL0}           & $C_L$               & lift at zero $\alpha$\\
\code{CLalpha}       & $C_L$               & lift-curve slope $\partial C_L/\partial\alpha$\\
\code{CLmax}         & $C_L$               & maximum lift coefficient\\
\code{alpha\_stall}  & $C_L$               & stall angle of attack\\
\code{CM0}           & $C_M$               & pitching moment at zero lift\\
\code{CMalpha}       & $C_M$               & pitch stiffness $\partial C_M/\partial\alpha$\\
\code{CD0}           & $C_D$, $C_L$        & zero-lift drag\\
\code{k}             & $C_D$, $C_L$        & induced-drag factor in $C_D = C_{D0} + k\,C_L^2$\\
\code{CL\_trim}      & $C_L$, $C_M$        & $C_L$ at $C_M = 0$\\
\bottomrule
\end{tabularx}

All extracted parameters carry propagated uncertainty when the input has
an UncFrame.
\end{reqbox}

\section{Pipeline Persistence}

\begin{reqbox}{REQ-53}{\stable}
\textbf{Indexed history: export by range.}\\
\code{db.history.to\_pipeline(start=None, end=None)} extracts a contiguous
range of history entries (by index) as a reusable Pipeline object.
Defaults to the full history. Examples:
\code{db.history.to\_pipeline()},
\code{db.history.to\_pipeline(start=2, end=7)},
\code{db.history.to\_pipeline(start=10)}.
\end{reqbox}

\begin{reqbox}{REQ-54}{\stable}
\textbf{Apply pipeline to a new VarFrame.}\\
\code{pipeline.apply(db\_new)} applies the recorded sequence to a new
VarFrame and returns the result. Missing variables raise
\code{PipelineCompatibilityError}.
\end{reqbox}

\begin{reqbox}{REQ-55}{\stable}
\textbf{Pipeline serialization.}\\
\code{pipeline.save(path)} writes a \code{.itc\_pipe} file.
\code{itc.load\_pipeline(path)} reads it. The file is human-readable and
version-controllable.
\end{reqbox}

\section{Plotting: \code{itc.datavis()}}

\begin{reqbox}{REQ-56}{\stable}
\textbf{DataVis object.}\\
\code{p = itc.datavis(backend="matplotlib", colormap=None, markercycle=None,
linecycle=None, figsize=None)} returns a DataVis object holding
figure-level style settings. \code{backend} is \code{"matplotlib"}
(default) or \code{"plotly"} (interactive).
\end{reqbox}

\begin{reqbox}{REQ-57}{\stable}
\textbf{Single-plot semantic axis mapping.}\\
\code{p.plot(db, x=, y=, DataAxis=, ColorAxis=, SymbolAxis=,
LineStyleAxis=, colorbar=False, ...)} iterates over all unique values of
\code{DataAxis}, plotting one curve per value. The \code{x} variable may
co-vary with \code{DataAxis}. Multiple calls overlay curves on the same
figure. Axis units (from \code{Dimension.unit} and \code{Variable.unit})
are appended to labels by default.
\end{reqbox}

\begin{reqbox}{REQ-58}{\stable}
\textbf{Multi-plot capability.}\\
\code{p.plot(..., multiplot=True, n\_cols=N)} produces a grid of subplots
when any of \code{db}, \code{x}, \code{y}, or \code{NFigAxis} is given as
a list:
\begin{itemize}[noitemsep,topsep=2pt]
  \item \code{db=[db1, db2, ...]}: one subplot per VarFrame;
  \item \code{x=[var1, var2, ...]}: one subplot per $x$ variable;
  \item \code{y=[var1, var2, ...]}: one subplot per $y$ variable;
  \item \code{NFigAxis="dim"}: one subplot per unique value of the
    named dimension.
\end{itemize}
The number of rows is computed as
$n_\text{rows} = \lceil n_\text{total} / n_\text{cols}\rceil$; trailing
empty cells are left blank. Default subplot size is
$3.5\,\text{in} \times 3\,\text{in}$.
\end{reqbox}

\begin{reqbox}{REQ-59}{\stable}
\textbf{Origin-tag visualization.}\\
When the input VarFrame has a HistoryFrame, DataVis distinguishes the
origin of plotted points by marker style: filled marker for \code{0}
(original), hollow marker for \code{+1} (interpolated/computed), hollow
marker with cross overlay for \code{-1} (extrapolated). Controlled by
\code{tag\_markers=True/False} (default \code{True}).
\end{reqbox}

\begin{reqbox}{REQ-60}{\stable}
\textbf{ItcFigure return object.}\\
\code{p.get\_figure()} returns an \code{ItcFigure} with
\code{.export(path, dpi=)}, \code{.add\_legend(loc=)},
\code{.annotate()}, and \code{.mpl} (escape hatch to the underlying
\code{matplotlib.Figure} or \code{plotly.graph\_objects.Figure}).
\end{reqbox}

\begin{reqbox}{REQ-61}{\stable}
\textbf{AIAA default style.}\\
The default style follows AIAA publication standards: serif fonts,
visible grid, white background, inward tick marks, mathtext labels (no
LaTeX dependency at plot time), figsize $7\,\text{in}\times 5\,\text{in}$
for single plots.
\end{reqbox}

\begin{reqbox}{REQ-62}{\stable}
\textbf{Specialized plot methods.}\\
DataVis exposes \code{plot\_polar()}, \code{plot\_cp()},
\code{plot\_disk()}, \code{plot\_colormap()}, and
\code{plot\_colormap\_overlay()}. All inherit the full option set of
\code{p.plot()}. Each method documents the AIAA convention it implements
(e.g.\ inverted $C_p$ axis).
\end{reqbox}

\section{Surrogate Modeling: \code{itc.surrogate()}}

\begin{reqbox}{REQ-63}{\stable}
\textbf{Surrogate factory.}\\
\code{varframe, model = itc.surrogate(db, inputs=[...], output=str,
model=str, options=\{\})} fits an SMT~\citep{smt} surrogate. The function
returns a tuple:
\begin{itemize}[noitemsep,topsep=2pt]
  \item the first element is the fitted VarFrame on the original grid
    (for inspection and comparison);
  \item the second element is the standalone \code{Surrogate} object,
    independently exportable.
\end{itemize}
Supported models: \code{"KRG"}, \code{"RBF"}, \code{"IDW"} at minimum.
\end{reqbox}

\begin{reqbox}{REQ-64}{\stable}
\textbf{Surrogate predict, validate, and persist.}\\
A fitted \code{Surrogate} supports:
\begin{itemize}[noitemsep,topsep=2pt]
  \item \code{.predict(**coords)} returning a VarFrame;
  \item \code{.cross\_validate(k=)} returning a results object with
    \code{.rmse} per fold;
  \item \code{.save(path, embed\_source=False)} writing a
    \code{.itc\_surr} file (optionally embedding the source VarFrame for
    self-contained audit);
  \item \code{itc.load\_surrogate(path)} restoring it.
\end{itemize}
\end{reqbox}

\section{Aerospace Framework: \code{itc.aerospace}}

\begin{reqbox}{REQ-65}{\draft}
\textbf{Aerospace subpackage.}\\
\code{itc.aerospace} is a subpackage containing engineering analysis
modules for aeronautical applications. Unlike \code{pproc} processors
(which process existing data), \code{itc.aerospace} modules implement
physical models that generate data in VarFrame format. Every result is a
VarFrame with full Provenance and History, making the entire computation
auditable.
\end{reqbox}

\begin{reqbox}{REQ-66}{\draft}
\textbf{Aerospace: performance.}\\
\code{itc.aerospace.performance} provides \code{TakeoffAnalysis},
\code{ClimbAnalysis}, \code{EnvelopeAnalysis}. Each accepts aerodynamic
and propulsion VarFrame inputs and produces a VarFrame output where every
timestep or flight condition is a datapoint. The complete computation
memory (forces, accelerations, velocities, distances) is auditable.
\end{reqbox}

\begin{reqbox}{REQ-67}{\draft}
\textbf{Aerospace: stabcontrol and EUB.}\\
\code{itc.aerospace.stabcontrol} provides static margin computation, trim
horizontal tail deflection with uncertainty propagation (the EUB on trim
deflection), pitch-damping derivatives, and V-n diagram generation.
Uncertainty in wind tunnel inputs (including correlations from a common
calibration) propagates automatically to uncertainty on required trim
deflections.
\end{reqbox}

\begin{reqbox}{REQ-68}{\draft}
\textbf{Aerospace: propeller BEMT.}\\
\code{itc.aerospace.propeller.BEMTAnalysis} implements Blade Element
Momentum Theory following Adkins \& Liebeck~\citep{adkinsLiebeck} with
Prandtl tip-loss correction. Inputs are blade geometry and airfoil polar
VarFrame objects. Output is a VarFrame containing radial distributions of
thrust and torque loading, integrated performance coefficients ($C_T$,
$C_P$, $\eta$), and (when manufacturing tolerances are assigned as uncertainties on blade geometry parameters) propagated uncertainty on
all outputs.
\end{reqbox}

\begin{reqbox}{REQ-69}{\draft}
\textbf{Auditable computation memory.}\\
All \code{itc.aerospace} analyses store the full computation memory in
the output VarFrame. For time-marching analyses, every timestep is a
datapoint in a time dimension. For iterative analyses (e.g.\ BEMT), every
radial station and iteration is accessible. The goal is that any output
number can be traced back to the exact input values and model steps that
produced it.
\end{reqbox}

\section{Output and Export}

\begin{reqbox}{REQ-70}{\stable}
\textbf{Export formats.}\\
A VarFrame supports export to:
\begin{itemize}[noitemsep,topsep=2pt]
  \item flat CSV (\code{to\_csv});
  \item JSON with optional uncertainty (\code{to\_json});
  \item pandas DataFrame (\code{to\_pandas});
  \item NumPy array with dimension labels
    (\code{to\_numpy(return\_dims=True)});
  \item native \itc{} binary \code{.itc} (\code{save} / \code{itc.open});
  \item PROV-N and PROV-JSON for the provenance graph
    (\code{db.export\_provenance(path, format="prov-json")}).
\end{itemize}
The \code{.itc} format preserves all metadata, Provenance, History,
uncertainty, correlation, and origin tags.
\end{reqbox}

\begin{reqbox}{REQ-71}{\stable}
\textbf{Provenance and history in all exports.}\\
Every export embeds metadata derived from Provenance and a human-readable
summary of History in the output file:
\begin{itemize}[noitemsep,topsep=2pt]
  \item CSV: header comment lines (\code{\# ITACA export: version: ...});
  \item JSON: top-level \code{"provenance"} and \code{"history"} keys;
  \item \code{.itc}: native JSON files inside the archive.
\end{itemize}
\end{reqbox}

\begin{reqbox}{REQ-72}{\stable}
\textbf{Split export by dimension.}\\
\code{db.to\_csv(dir, split\_by=dim)} writes one CSV file per unique
value of the specified dimension.
\end{reqbox}

\begin{reqbox}{REQ-73}{\stable}
\textbf{Unit conversion utility.}\\
\code{utils.units.convert(value, from\_unit, to\_unit)} converts numeric
values between units using \itc{}'s internal conversion table. The table
is hand-curated, validated against unit tests, and covers SI base units
plus common aerospace units (deg, rpm, knot, ft, lb, slug, etc.).
External unit libraries are not used (DD-13, NREQ-09).
\end{reqbox}

\section{Library Infrastructure}
\label{sec:library-infrastructure}

The requirements of this section are adoptions from the 2026-07-23
library-review triage (decisions D1, D2, and D6; DD-24 records the
options-registry adoption): infrastructure patterns proven in the
scientific Python ecosystem, specified here so that they enter \itc{}
through the same requirement process as every domain feature.

\begin{reqbox}{REQ-104}{\draft}
\textbf{Options registry with exact keys.}\\
\itc{} provides a central options registry:
\code{itc.set\_option(key, value)}, \code{itc.get\_option(key)},
\code{itc.describe\_option(key)}, and the context manager
\code{itc.option\_context(key, value, ...)}. Keys are exact,
dot-namespaced strings (for example \code{"plot.theme"}); there is no
partial or abbreviated matching, and an unknown key raises
\code{OptionsError} naming the key, the closest registered keys, and
\code{describe\_option} as the fix. Every registered option declares
its type, its validator, and its default at registration; a validator
rejection names the offered value, the accepted domain, and the
bounds in the message. An option may be registered as
required-until-set: reading it before it is set raises with the owner prefix and the
setting call in the message. The registry supports snapshot and
restore for tests; the test suite resets it through an autouse
fixture. Registration of a duplicate key raises.
\end{reqbox}

\begin{reqbox}{REQ-105}{\draft}
\textbf{Typed no-default sentinel.}\\
\itc{} defines a single sentinel object \code{itc.no\_default} whose
type is expressible in annotations (an enum singleton usable with
\code{Literal}), used by every public signature that must distinguish
an argument that was not passed from an argument explicitly set to
\code{None}. The sentinel's repr is readable
(\code{<no\_default>}); it is never a valid data value; passing it
explicitly where it is not the declared default raises. All M1
operations adopt it where the distinction changes behavior, and new
public signatures adopt it from day one.
\end{reqbox}

\begin{reqbox}{REQ-106}{\draft}
\textbf{Accessor registration.}\\
\code{itc.register\_accessor(name)} is a class decorator registering
a named accessor namespace on VarFrame: after registration,
\code{db.<name>} instantiates the accessor class with the frame and
caches it per instance. Registering a name that collides with an
existing attribute or accessor raises at registration time, naming
the holder. An \code{AttributeError} escaping the accessor's
\code{\_\_init\_\_} is re-raised as \code{RuntimeError} carrying the
original, so attribute machinery never swallows a real defect
silently. Accessors are the sanctioned extension point for external
packages; the sister-library page documents the contract. The
registry supports snapshot and restore for tests.
\end{reqbox}
